\documentclass[noanswers, addpoints, 12pt]{exam}
\usepackage[T1]{fontenc}
\usepackage[utf8]{inputenc}
\usepackage{lmodern}
\usepackage{pgfplots}
\pgfplotsset{compat=1.17}
\usepackage{tikz}
\usetikzlibrary{positioning,shapes.multipart,arrows.meta,shadows}
\usepackage[french]{babel}
\usepackage{xcolor}
\usepackage{listings}
\usepackage{graphicx}
\usepackage[most]{tcolorbox}

\definecolor{codegreen}{rgb}{0,0.6,0}
\definecolor{codegray}{rgb}{0.5,0.5,0.5}
\definecolor{codepurple}{rgb}{0.58,0,0.82}
\definecolor{backcolour}{rgb}{0.95,0.95,0.92}
 
\lstdefinestyle{mystyle}{
    backgroundcolor=\color{backcolour},   
    commentstyle=\color{codegreen},
    keywordstyle=\color{magenta},
    numberstyle=\tiny\color{codegray},
    stringstyle=\color{codepurple},
    basicstyle=\footnotesize,
    breakatwhitespace=false,         
    breaklines=true,                 
    captionpos=b,                    
    keepspaces=true,                 
    numbers=left,                    
    numbersep=5pt,                  
    showspaces=false,                
    showstringspaces=false,
    showtabs=false,                  
    tabsize=2
}
\lstset{style=mystyle, language=Python}

\title{CC2 Term NSI - SQL, Tri et complexité}
\author{Babash - NSI Tani Malandi et Johanne Bakalara}
\date{Janvier 2024, Durée: 1h}


\begin{document}
\runningfooter{}{\footnotesize Babash - NSI Tani Malandi et Johanne Bakalara \textperiodcentered{} CC BY-SA 4.0 \textperiodcentered{} page \thepage}{}
\firstpagefooter{}{\footnotesize Babash - NSI Tani Malandi et Johanne Bakalara \textperiodcentered{} CC BY-SA 4.0 \textperiodcentered{} page \thepage}{}

\maketitle

\pagestyle{headandfoot}
\thispagestyle{headandfoot}
\firstpageheader{Nom : \\  Prénom :\\ Classe :}{}{}


\section*{Consignes:}
\begin{center} 
\fbox{\parbox{\linewidth}{
\begin{center}
Le devoir est à rédiger sur une feuille libre.\\
La clarté de la rédaction et les explications font partie intégrante de la notation. Une réponse non rédigée sans explication sera considérée fausse.\\
Prenez le temps d'encadrer les réponses.
\end{center}


}}
\end{center}

\vsword{Note}
\hsword{Note}

\section*{Barème :}
\begin{center}
  \par\bigskip
  \gradetable[h]
\end{center}
\begin{questions}

\section{SQL : }

\question[12] Une bibliothèque gère une base de données dont on donne les tables et les attributs :\\
    \begin{itemize}
        \item \indent Abonne(\underline{Abonne\_ID}, Nom, Prenom, DateNaissance, Adresse, Ville, CodePostal)
        \item \indent Livre(\underline{Livre\_ID}, Titre, Auteur, Genre, ISBN)
        \item \indent Emprunt(\#Abonne\_ID, \#Livre\_ID, DateEmprunt, DateRetourPrevu, DateRetourReel)
    \end{itemize}
\begin{center}
\begin{center}
\begin{tikzpicture}[font=\footnotesize,tab/.style={draw,rectangle split,rectangle split parts=2,align=left,inner sep=3pt}]
\node[tab] (a) at (0,0) {\textbf{Abonne}\nodepart{second}\textbf{Abonne\_ID} : INT PRIMARY KEY\\Nom : VARCHAR(50)\\Prenom : VARCHAR(50)\\DateNaissance : DATE\\Adresse : VARCHAR(255)\\Ville : VARCHAR(50)\\CodePostal : CHAR(5)};
\node[tab] (l) at (8,0) {\textbf{Livre}\nodepart{second}\textbf{Livre\_ID} : INT PRIMARY KEY\\Titre : VARCHAR(100)\\Auteur : VARCHAR(100)\\Genre : VARCHAR(255)\\ISBN : VARCHAR(18)};
\node[tab] (e) at (4,-4.2) {\textbf{Emprunt}\nodepart{second}\textbf{\#Abonne\_ID} : INT\\\textbf{\#Livre\_ID} : INT\\DateEmprunt : DATE\\DateRetourPrevu : DATE\\DateRetourReel : DATE};
\draw[->,thick] (e.west) -| (a.south);
\draw[->,thick] (e.east) -| (l.south);
\end{tikzpicture}
\end{center}

\end{center}
Le format des dates est YYYY-MM-DD, par exemple '2020-11-02' pour le 2 novembre 2020.


\begin{parts}

\subsection{Sélection}

\part Sélectionner les noms et prénoms des abonnés habitant à Chirongui ?

\part Sélectionner le nom des gens dont le code postal commence par 976 ?

\part Sélectionner le nom des gens ayant emprunté le livre Papillon ?

\part Sélectionner tous les attributs de l'abonné le plus âgé ?



\subsection{Modification}
\part Quelle requête permet de rajouter l’emprunt du livre d’identifiant 4252 par l'abonné d'identifiant 803 le 2 novembre 2020 sachant que les emprunts sont pour une durée de trois semaines ?

\subsection{Ajout d'une table}
\part On souhaite rajouter une table Emplacement permettant de représenter l'emplacement du livre dans la bibliothèque, la table aura les attributs suivants :
\begin{description}
    \item[id]: Le numéro du livre (INT)
    \item[zone]: le nom de la zone (VARCHAR 50)
    \item[etagere]: le numéro d'étagère (INT)
\end{description}


\end{parts}

\newpage

\section{Algorithme de tri :}

\begin{lstlisting}[language=Python]
def tri_surprise(tab):
    if len(tab) < 2:
        return tab
    for i in range(1, len(tab)):
        e = tab[i]
        j = i - 1
        while j >= 0 and e < tab[j] :
            tab[j+1] = tab[j]
            j = j - 1
        tab[j+1] = e
    return tab
\end{lstlisting}

\question[2] Une seule réponse possible :
\begin{parts}
    \part Quel est cet algorithme de tri ?
    \begin{choices}
        \choice Tri par sélection
        \correctchoice Tri par insertion
        \choice Tri fusion
    \end{choices}
    \part Quelle est la complexité de cet algorithme ?
    \begin{choices}
        \choice Logarithmique $O(log(n))$
        \choice Linéaire $O(n)$
        \choice Quasi-linéaire $O(nlog(n))$
        \correctchoice Quadratique $O(n^2)$
        \choice Exponentiel $O(2^n)$
    \end{choices}
\end{parts}

\question[3] Associer les équations suivantes aux courbes correspondantes :\\
$f(x)=10$, $f(x)=x$, $f(x)=\log_2(x)$, $f(x)=x\log_2(x)$, $f(x)=x^2$, $f(x)=2^x$. Chaque courbe est repérée par une lettre.\\
\begin{center}
\begin{tikzpicture}
\begin{axis}[width=0.95\linewidth,height=8cm,xmin=0,xmax=20,ymin=0,ymax=40,xlabel={$x$},ylabel={$f(x)$},grid=both,minor tick num=1,
 every axis plot/.append style={thick},domain=0.01:20,samples=200,legend pos=north west,legend style={font=\footnotesize}]
\addplot[color=red] {10};
\addplot[color=blue] {x};
\addplot[color=green!50!black] {ln(x)/ln(2)};
\addplot[color=gray] {x*ln(x)/ln(2)};
\addplot[color=violet,domain=0:7] {x^2};
\addplot[color=orange,domain=0:6] {2^x};
\legend{A,B,C,D,E,F}
\end{axis}
\end{tikzpicture}
\end{center}


\question[2] Donner la valeur de tab à chaque itération de la boucle while pour l'appel suivant : tri\_surprise([1, 5, 100, 78, 8, 22])


\question[3] Compléter la fonction fusion ci-dessous.
\begin{lstlisting}[language=Python]
def fusion(gauche, droite, tab):
    i = ...
    j = ...
    while i + j < len(tab):
        if j == len(droite) or (i < len(gauche) and gauche[i] < droite[j]):
            tab[i + j] = ...
            i = i + 1
        else:
            tab[i + j] = droite[j]
            j = ...

def tri_fusion(tab):
    if len(tab) < 2:
        return tab
    milieu = len(tab) // 2
    gauche = tab[:milieu]
    droite = tab[milieu:]
    tri_fusion(gauche)
    tri_fusion(droite)
    fusion(gauche, droite, tab)

\end{lstlisting}

\end{questions}

\end{document}

